\chapter{Redis and Caching Topologies}
\label{ch:redis-caching}

\begin{chaptergoalsbox}
\textbf{What you will learn:}
\begin{itemize}[leftmargin=12pt]
  \item Why Redis uses a single-threaded non-blocking I/O event loop (\texttt{epoll}/\texttt{kqueue}) and how its C structures (\texttt{sds}, \texttt{dict}, \texttt{skiplist}, and \texttt{listpack} replacing \texttt{ziplist}) optimize RAM and CPU cache locality.
  \item Approximated LRU and probabilistic Morris-counter LFU memory eviction.
  \item Caching topologies in Python: Cache-Aside, Read-Through, Write-Through, and Write-Behind.
  \item Defending against Stampedes (XFetch), Penetration (Bloom filters), Avalanches (TTL jitter), and atomic Lua scripts using \texttt{HSET}.
  \item Redis Sentinel failover vs. Redis Cluster 16,384 hash slots, hash tags, and \texttt{-MOVED} redirects.
\end{itemize}

\textbf{What you should already understand:}
\begin{itemize}[leftmargin=12pt]
  \item Basic key-value operations (\texttt{GET}, \texttt{SET}, \texttt{DEL}) and relational SQL queries.
  \item Network round-trip costs from Chapter~\ref{ch:distributed-fundamentals}; no prior C data-structure internals are assumed.
\end{itemize}

\textbf{Why this topic exists:}
Reading a row from DRAM takes $\sim 100\text{ ns}$, whereas querying a relational database on disk or joining tables under lock contention takes milliseconds---ten thousand times slower. Yet teams that slap Redis in front of a database without understanding eviction, stampedes, or write invalidation routinely trade slow pages for silent stale data and midnight outages.

\textbf{Motivating practical problem:}
An airline ticketing portal caches flight-availability searches with a fixed 24-hour TTL expiring at 00:00 UTC. At 00:00:01, 80,000 concurrent browsers miss Redis simultaneously and slam PostgreSQL with identical complex queries. Database connections saturate, disk queues spike, and the booking platform stays down for 45 minutes.
\end{chaptergoalsbox}

\section{Redis Architecture and Internal Data Structures}

Created by Salvatore Sanfilippo in 2009, \textbf{Redis} (Remote Dictionary Server) stores typed data structures directly in RAM while optionally persisting snapshots or append logs to disk.

\subsection{The Single-Threaded Multiplexed Event Loop}

Contrary to intuition, Redis processes all client commands on a \textbf{single main thread} (while offloading slow background tasks like large-key deletion via \texttt{UNLINK} and disk \texttt{fsync} to helper threads). Instead of spawning one OS thread per client socket, Redis registers thousands of non-blocking TCP sockets with the kernel's \textbf{I/O multiplexer} (\texttt{epoll} on Linux, \texttt{kqueue} on BSD/macOS), waking up only for sockets that actually have bytes ready to read or write.

\begin{figure}[htbp]
\centering
\begin{tikzpicture}[
  scale=0.85,
  sock/.style={rectangle, draw=headingblue, fill=codebg, thick, rounded corners=2pt, minimum width=2.4cm, minimum height=0.7cm, align=center, font=\sffamily\tiny},
  loop/.style={circle, draw=exampleBlue, fill=exampleBg, thick, minimum size=3.0cm, align=center, font=\sffamily\footnotesize}
]
  \node[sock] (s1) at (-4, 1.2) {Client Socket 1};
  \node[sock] (s2) at (-4, 0.0) {Client Socket 2};
  \node[sock] (s3) at (-4, -1.2) {Client Socket 3};

  \node[rectangle, draw=hoodSlate, fill=hoodBg, thick, rounded corners=2pt, minimum width=2.3cm, minimum height=3.2cm, align=center, font=\sffamily\footnotesize] (mux) at (-0.5, 0) {I/O Multiplexer\\(\texttt{epoll} / \texttt{kqueue})};

  \node[loop] (el) at (4.5, 0) {\textbf{Event Loop}\\Sequential Exec\\In-Memory Dict\\(Zero Locks)};

  \draw[-{Stealth[scale=1.0]}, thick] (s1) -- (mux);
  \draw[-{Stealth[scale=1.0]}, thick] (s2) -- (mux);
  \draw[-{Stealth[scale=1.0]}, thick] (s3) -- (mux);

  \draw[-{Stealth[scale=1.0]}, thick, color=exampleBlue] (mux) -- node[above, font=\sffamily\tiny] {Ready Events} (el);
\end{tikzpicture}
\caption{Redis Single-Threaded Non-Blocking Event Loop}
\label{fig:v4-ch06-event-loop}
\end{figure}

Why is single-threaded execution so fast? In-memory operations take nanoseconds, so Redis is bottlenecked by network I/O and RAM bandwidth, not CPU cores. Running commands sequentially eliminates mutex locks, deadlocks, and OS thread context-switching overhead, letting one core sustain $>100{,}000$ operations per second.

\subsection{Internal C Data Structures: SDS, Dict, SkipList, and Listpack}

\begin{itemize}[leftmargin=18pt]
  \item \textbf{SDS (Simple Dynamic String)}: In standard C, strings end at the first null byte (\texttt{'\textbackslash 0'}), forcing $O(N)$ scans to compute length and breaking on binary payloads. Redis wraps strings in \texttt{sdshdr}, storing exact length and allocated capacity upfront---giving $O(1)$ length checks, binary safety (images, Protobufs), and pre-allocated growth headroom.
  \item \textbf{Dict (Incremental Rehashing Hash Table)}: When a hash table fills up and must double in size, moving 50 million keys at once would freeze the single-threaded event loop for seconds. Redis keeps \emph{two} hash tables ($ht[0]$ and $ht[1]$) and performs \textbf{incremental rehashing}: each incoming client command migrates just one bucket from $ht[0]$ to $ht[1]$ until the resize completes smoothly.
  \item \textbf{SkipList (Sorted Sets / \texttt{ZSET})}: Imagine a highway with local exits at every mile (Level~1), express exits every few miles (Level~2), and interstate hubs (Level~3). Instead of complex tree rotations (like Red-Black trees), a SkipList builds these express lanes probabilistically: every inserted element starts at Level~1 and repeatedly ``flips a biased coin'' (probability $p = 0.25$ in Redis, up to 32 levels) to get promoted to higher express lanes. Searching from the top lane down achieves $O(\log N)$ lookups and fast range scans (\texttt{ZRANGE}).
  \item \textbf{Listpack (and Why \texttt{ziplist} Was Replaced)}: Small hashes and lists are packed contiguously into a single byte array for CPU cache efficiency. Historically, Redis used \texttt{ziplist}, where each entry stored the byte length of the \emph{previous} entry (\texttt{prevlen}) so Redis could walk backward. If \texttt{prevlen} was $< 254$ bytes, it took 1 byte; if $\ge 254$ bytes, it took 5 bytes. Inserting one 254-byte item near the front forced the next item's \texttt{prevlen} to grow from 1 to 5 bytes---pushing *that* item over 254 bytes and triggering an $O(N^2)$ \textbf{cascading reallocation} down the list! Redis~7 replaced \texttt{ziplist} with \textbf{\texttt{listpack}}, where every entry stores its \emph{own} encoding and total length at its tail. Walking backward reads the tail length in $O(1)$ with zero dependency on neighbors, completely eliminating cascading updates.
\end{itemize}

\begin{figure}[htbp]
\centering
\begin{tikzpicture}[
  scale=0.85,
  node/.style={rectangle, draw=headingblue, fill=codebg, thick, rounded corners=2pt, minimum width=0.9cm, minimum height=0.6cm, align=center, font=\sffamily\tiny}
]
  \node[node, fill=warningBg] (l3_0) at (0, 2) {Head};
  \node[node, fill=warningBg] (l3_4) at (6, 2) {40};
  \node[node, fill=warningBg] (l3_8) at (10, 2) {NIL};
  \draw[-{Stealth[scale=1.0]}, thick] (l3_0) -- (l3_4);
  \draw[-{Stealth[scale=1.0]}, thick] (l3_4) -- (l3_8);

  \node[node, fill=intuitionBg] (l2_0) at (0, 1) {Head};
  \node[node, fill=intuitionBg] (l2_2) at (3, 1) {20};
  \node[node, fill=intuitionBg] (l2_4) at (6, 1) {40};
  \node[node, fill=intuitionBg] (l2_7) at (8.5, 1) {70};
  \node[node, fill=intuitionBg] (l2_8) at (10, 1) {NIL};
  \draw[-{Stealth[scale=1.0]}, thick] (l2_0) -- (l2_2);
  \draw[-{Stealth[scale=1.0]}, thick] (l2_2) -- (l2_4);
  \draw[-{Stealth[scale=1.0]}, thick] (l2_4) -- (l2_7);
  \draw[-{Stealth[scale=1.0]}, thick] (l2_7) -- (l2_8);

  \node[node] (l1_0) at (0, 0) {Head};
  \node[node] (l1_1) at (1.5, 0) {10};
  \node[node] (l1_2) at (3, 0) {20};
  \node[node] (l1_3) at (4.5, 0) {30};
  \node[node] (l1_4) at (6, 0) {40};
  \node[node] (l1_5) at (7.2, 0) {50};
  \node[node] (l1_7) at (8.5, 0) {70};
  \node[node] (l1_8) at (10, 0) {NIL};
  \draw[-{Stealth[scale=1.0]}, thick] (l1_0) -- (l1_1);
  \draw[-{Stealth[scale=1.0]}, thick] (l1_1) -- (l1_2);
  \draw[-{Stealth[scale=1.0]}, thick] (l1_2) -- (l1_3);
  \draw[-{Stealth[scale=1.0]}, thick] (l1_3) -- (l1_4);
  \draw[-{Stealth[scale=1.0]}, thick] (l1_4) -- (l1_5);
  \draw[-{Stealth[scale=1.0]}, thick] (l1_5) -- (l1_7);
  \draw[-{Stealth[scale=1.0]}, thick] (l1_7) -- (l1_8);
\end{tikzpicture}
\caption{Multi-Level SkipList Structure Supporting $O(\log N)$ Sorted Set Lookups}
\label{fig:v4-ch06-skiplist}
\end{figure}

\section{Memory Eviction: Approximated LRU and Morris LFU}

When Redis memory hits \texttt{maxmemory}, it applies the configured eviction policy (\texttt{noeviction}, \texttt{allkeys-lru}, \texttt{volatile-lru}, \texttt{allkeys-lfu}, or \texttt{volatile-lfu}).

\begin{underthehoodbox}[Sampled LRU and the Morris Counter Intuition for LFU]
True LRU needs a global doubly linked list updated on every read, wasting 16 bytes of pointers per key. Instead, Redis stores a 24-bit clock field inside each object and samples $N$ random keys (default $N=5$), evicting the oldest among the sample. With $N=10$, accuracy matches exact LRU within 1\%.

Under \textbf{LFU (Least Frequently Used)}, Redis packs both a 16-bit decay timestamp and an 8-bit frequency counter into that same 24-bit field. How can an 8-bit integer (max value $255$) distinguish $1{,}000$ hits from $1{,}000{,}000$ hits? Using a probabilistic \textbf{Morris counter}: when a key at counter $c$ is accessed, Redis increments $c$ only with probability $p = \frac{1}{(c - 5) \cdot \text{lfu-log-factor} + 1}$. As $c$ grows, increments become rarer---turning 8 bits into a compact logarithmic estimator of millions of requests!
\end{underthehoodbox}

\section{Caching Topologies with Runnable Code}

\begin{figure}[htbp]
\centering
\begin{tikzpicture}[
  scale=0.84,
  box/.style={rectangle, draw=headingblue, fill=codebg, thick, rounded corners=2pt, minimum width=2.1cm, minimum height=0.85cm, align=center, font=\sffamily\tiny},
  db/.style={cylinder, draw=exampleBlue, fill=exampleBg, thick, shape border rotate=90, minimum width=1.6cm, minimum height=1.1cm, align=center, font=\sffamily\tiny}
]
  % Left: Cache-Aside
  \node[anchor=west, font=\sffamily\bfseries\footnotesize, color=headingblue] at (-0.6, 2.3) {(a) Cache-Aside (App-Managed)};
  \node[box] (app1) at (0.3, 0.4) {Application};
  \node[box, fill=intuitionBg] (c1) at (3.6, 1.4) {Redis Cache};
  \node[db] (db1) at (3.6, -0.7) {Database};
  \draw[-{Stealth[scale=0.9]}, thick] (app1) -- node[above, sloped, font=\sffamily\tiny] {1. GET} (c1);
  \draw[-{Stealth[scale=0.9]}, thick, color=warningRed] (app1) -- node[below, sloped, font=\sffamily\tiny] {2. Miss $\to$ SQL} (db1);
  \draw[-{Stealth[scale=0.9]}, thick, dashed, color=intuitionTeal] (app1) to [bend left=35] node[above, font=\sffamily\tiny] {3. SET} (c1);

  % Right: Write-Through vs Write-Behind
  \node[anchor=west, font=\sffamily\bfseries\footnotesize, color=headingblue] at (6.4, 2.3) {(b) Write-Through / Write-Behind (Cache-Managed)};
  \node[box] (app2) at (7.3, 0.4) {Application};
  \node[box, fill=exampleBg] (c2) at (10.7, 0.4) {Cache Proxy /\\Provider};
  \node[db] (db2) at (14.2, 0.4) {Database};
  \draw[-{Stealth[scale=0.9]}, thick] (app2) -- node[above, font=\sffamily\tiny] {Read/Write} (c2);
  \draw[-{Stealth[scale=0.9]}, thick, color=exampleBlue] (c2) -- node[above, font=\sffamily\tiny] {Sync (Write-Through)} (db2);
  \draw[-{Stealth[scale=0.9]}, thick, dashed, color=mistakeAmber] (c2) to [bend right=30] node[below, font=\sffamily\tiny] {Async Batch (Write-Behind)} (db2);
\end{tikzpicture}
\caption{Cache-Aside vs. Write-Through and Write-Behind Topologies}
\label{fig:v4-ch06-write-topologies}
\end{figure}

\begin{enumerate}[leftmargin=18pt]
  \item \textbf{Cache-Aside (Lazy Loading)}: The application checks Redis first; on a miss, it queries PostgreSQL and writes back to Redis with a jittered TTL. On writes, it updates PostgreSQL and \emph{deletes} the cache key.
  \item \textbf{Read-Through}: The application talks only to a caching library/proxy, which automatically loads from the database on misses so caller code stays clean.
  \item \textbf{Write-Through}: Writes go through the cache, which updates the database synchronously before acknowledging---strong consistency, higher write latency.
  \item \textbf{Write-Behind (Write-Back)}: Writes return immediately after updating cache memory; a background flusher batches updates to the database asynchronously. Extreme write speed, but risks data loss if the cache node crashes before flushing.
\end{enumerate}

\begin{lstlisting}[language=Python, caption={Cache-Aside and Read-Through with TTL Jitter in Python}]
import json
import random
from typing import Callable, Optional

def get_flight_cache_aside(redis_client, db, route_key: str, base_ttl: int = 300) -> dict:
    cache_key = f"flight:{route_key}"
    cached = redis_client.get(cache_key)
    if cached is not None:
        return json.loads(cached)

    # Cache miss: query relational DB and store with TTL jitter
    row = db.fetch_flight(route_key)
    jitter = random.randint(0, 60)  # Spread expirations across 60s window
    redis_client.set(cache_key, json.dumps(row), ex=base_ttl + jitter)
    return row

def update_flight_price(redis_client, db, route_key: str, new_price: int) -> None:
    db.update_price(route_key, new_price)
    redis_client.delete(f"flight:{route_key}")  # Invalidate, never SET blindly!

class ReadThroughCache:
    def __init__(self, redis_client, loader: Callable[[str], Optional[dict]], ttl: int = 300):
        self.redis = redis_client
        self.loader = loader
        self.ttl = ttl

    def get(self, key: str) -> Optional[dict]:
        raw = self.redis.get(key)
        if raw is not None:
            return json.loads(raw)
        data = self.loader(key)
        if data is not None:
            self.redis.set(key, json.dumps(data), ex=self.ttl + random.randint(0, 30))
        return data
\end{lstlisting}

\textbf{Step-by-step trace:}
\begin{enumerate}[leftmargin=18pt]
  \item Call \texttt{get\_flight\_cache\_aside(..., "NYC-LON")}: checks \texttt{flight:NYC-LON}. On a cold cache, \texttt{cached} is \texttt{None}, so it queries \texttt{db.fetch\_flight("NYC-LON")} (returning \texttt{\{"price": 420\}}), picks random \texttt{jitter = 23}, stores the JSON with \texttt{ex=323} seconds, and returns the dict.
  \item Subsequent calls hit Redis directly until 323s elapse. Because 10,000 keys receive random TTLs between 300s and 360s instead of expiring together, a \textbf{Cache Avalanche} at midnight is prevented.
  \item Calling \texttt{update\_flight\_price} commits to PostgreSQL and \emph{deletes} the cache key rather than calling \texttt{SET}, avoiding the stale-overwrite race where an older thread writes an outdated value after a newer thread updated the DB.
\end{enumerate}

\section{Cache Pathologies and Defenses (XFetch, Bloom Filters, Lua)}

\begin{table}[htbp]
\centering\small
\begin{tabular}{lp{4.8cm}p{6.2cm}}
\toprule
\textbf{Pathology} & \textbf{Mechanism} & \textbf{Production Defense} \\
\midrule
\textbf{Cache Stampede (Thundering Herd)} & Hot key expires; thousands of concurrent threads miss together and hit DB. & \textbf{XFetch} probabilistic early refresh or singleflight/mutex lock. \\
\textbf{Cache Penetration} & Attackers query nonexistent keys (e.g., ID \texttt{-999}) missing both cache and DB. & \textbf{Bloom Filter} ingress guard or caching short-TTL null tombstones. \\
\textbf{Cache Avalanche} & Huge batch of keys expires at the same timestamp (e.g., 00:00 UTC). & \textbf{TTL Jitter}: $\text{TTL} = \text{base} + \text{rand}(0, \Delta)$. \\
\textbf{Cache Breakdown} & One extreme celebrity hot key expires under huge QPS. & Async background refresh before expiry or logical soft-TTL inside value. \\
\bottomrule
\end{tabular}
\caption{Taxonomy of Distributed Caching Pathologies and Defenses}
\label{tab:v4-ch06-pathologies}
\end{table}

\subsection{Probabilistic Early Expiration: The XFetch Algorithm}

Why does $-\beta \cdot \delta \cdot \ln(\text{rand}()) > \text{TTL}_{\text{remaining}}$ stop stampedes without locks? Because when $u \sim \text{Uniform}(0, 1]$, $-\ln(u)$ follows an exponential distribution (mostly near $0.5\text{--}2.0$, occasionally larger). Multiplying by $\beta \cdot \delta$ (where $\delta$ is how long the DB query takes to recompute) means: when $\text{TTL}_{\text{remaining}}$ is large (say $250\text{ s}$), the threshold is virtually never crossed; as $\text{TTL}_{\text{remaining}}$ shrinks down toward $\delta$ (say $0.2\text{ s}$), \emph{one} lucky reader triggers an early refresh and resets the TTL before the key actually expires for everyone else!

\begin{figure}[htbp]
\centering
\begin{tikzpicture}[
  scale=0.86,
  axis/.style={->, thick, color=headingblue}
]
  \draw[axis] (0, 0) -- (11.5, 0) node[right, font=\sffamily\tiny] {Time};
  \draw[axis] (0, 0) -- (0, 2.4) node[above, font=\sffamily\tiny] {DB QPS};

  \draw[thick, color=warningRed] (4.5, 0) -- (4.7, 2.1) -- (5.0, 0);
  \node[font=\sffamily\tiny, color=warningRed, align=center] at (4.8, 2.3) {Without XFetch:\\All readers miss at $t=\text{TTL}$};

  \draw[thick, color=intuitionTeal] (3.4, 0) -- (3.4, 0.5) -- (3.5, 0);
  \node[font=\sffamily\tiny, color=intuitionTeal, align=center] at (1.9, 1.0) {XFetch:\\1 early refresh at $t=3.4$\\resets TTL smoothly};
  \draw[dashed, color=gray!60] (4.7, -0.2) -- (4.7, 1.9);
\end{tikzpicture}
\caption{Cache Stampede Spike vs. Smooth XFetch Early Refresh}
\label{fig:v4-ch06-stampede-xfetch}
\end{figure}

\begin{lstlisting}[language=Python, caption={XFetch Stampede-Resistant Reader in Python}]
import json, math, random, time

def get_with_xfetch(redis_client, db, key: str, ttl: float = 60.0, beta: float = 1.0):
    raw = redis_client.get(key)
    now = time.time()
    if raw is not None:
        entry = json.loads(raw)
        ttl_rem = entry["expiry"] - now
        delta = entry["delta"]
        u = max(random.random(), 1e-12)  # Avoid log(0)
        if ttl_rem > 0 and (-beta * delta * math.log(u)) <= ttl_rem:
            return entry["value"]

    # Cache miss OR probabilistically selected for early refresh
    t0 = time.monotonic()
    value = db.compute_expensive_report(key)
    delta = time.monotonic() - t0
    payload = {"value": value, "delta": delta, "expiry": time.time() + ttl}
    redis_client.set(key, json.dumps(payload), ex=int(ttl) + 5)
    return value
\end{lstlisting}

\textbf{Step-by-step trace:} Suppose \texttt{ttl = 60.0}, \texttt{beta = 1.0}, and \texttt{entry = \{"value": 99, "delta": 0.5, "expiry": now + 0.3\}} (so \texttt{ttl\_rem = 0.3s}). If a reader draws \texttt{u = 0.4}, then $-\ln(0.4) \approx 0.916$, and $1.0 \times 0.5 \times 0.916 = 0.458 > 0.3$. That single reader falls through to lines~15--19, recomputes \texttt{value}, measures \texttt{delta}, and resets \texttt{expiry = time.time() + 60.0} while all other concurrent readers still get the cached \texttt{99} cleanly.

\subsection{Bloom Filter Ingress Guard}

A \textbf{Bloom filter} is a compact bit array of $m$ bits paired with $k$ independent hash functions. It supports two answers: ``\emph{definitely not in the set}'' (100\% certain, blocking nonexistent-ID attacks before they reach PostgreSQL) or ``\emph{probably in the set}'' (small tunable false-positive rate $p \approx (1 - e^{-kn/m})^k$).

\begin{lstlisting}[language=Python, caption={Tiny Concrete Bloom Filter Ingress Guard in Python}]
import hashlib

class TinyBloomFilter:
    def __init__(self, size: int = 16, num_hashes: int = 2):
        self.size = size
        self.num_hashes = num_hashes
        self.bits = [0] * size

    def _positions(self, item: str) -> list[int]:
        pos = []
        for seed in range(self.num_hashes):
            digest = hashlib.sha256(f"{seed}:{item}".encode("utf-8")).hexdigest()
            pos.append(int(digest, 16) % self.size)
        return pos

    def add(self, item: str) -> None:
        for idx in self._positions(item):
            self.bits[idx] = 1

    def might_contain(self, item: str) -> bool:
        return all(self.bits[idx] == 1 for idx in self._positions(item))
\end{lstlisting}

\textbf{Step-by-step trace:} Create \texttt{bf = TinyBloomFilter(size=16, num\_hashes=2)}. Suppose \texttt{bf.\_positions("user:101")} hashes to indices \texttt{[3, 11]}. Calling \texttt{bf.add("user:101")} sets \texttt{bits[3] = 1} and \texttt{bits[11] = 1}. Later, an attacker sends \texttt{"user:-999"}, whose positions hash to \texttt{[3, 8]}. Since \texttt{bits[8] == 0}, \texttt{might\_contain("user:-999")} returns \texttt{False} immediately---rejecting the request in microseconds without touching Redis or PostgreSQL!

\subsection{Atomic Lua Scripts in Redis}

Because Redis runs one command at a time, sending \texttt{GET} then \texttt{SET} from Python allows another client to interleave in between. Sending a **Lua script** runs all steps atomically inside the single-threaded event loop. Note that Redis~4.0+ deprecated \texttt{HMSET} in favor of multi-field **`HSET`**.

\begin{lstlisting}[language=Python, caption={Atomic Lua Cache-Counter / Token Bucket Update Using \texttt{HSET}}]
LUA_TOKEN_BUCKET = """
local key = KEYS[1]
local capacity = tonumber(ARGV[1])
local rate = tonumber(ARGV[2])
local now = tonumber(ARGV[3])

local vals = redis.call('HMGET', key, 'tokens', 'ts')
local tokens = tonumber(vals[1]) or capacity
local last_ts = tonumber(vals[2]) or now
tokens = math.min(capacity, tokens + math.max(0, now - last_ts) * rate)

if tokens >= 1 then
    tokens = tokens - 1
    redis.call('HSET', key, 'tokens', tokens, 'ts', now)
    redis.call('EXPIRE', key, math.ceil(capacity / rate))
    return 1
end
return 0
"""
\end{lstlisting}

\textbf{Step-by-step trace:} Passing \texttt{KEYS[1]="rl:api:k1"}, \texttt{ARGV=["10", "2", "1000.0"]}: on the first call, \texttt{HMGET} returns \texttt{[nil, nil]}, initializing \texttt{tokens=10, last\_ts=1000.0}. Since \texttt{tokens >= 1}, it decrements \texttt{tokens} to \texttt{9}, atomically writes both fields with \texttt{HSET}, sets a 5-second TTL, and returns \texttt{1}.

\section{Redis Clustering: Sentinel vs. Cluster}

\textbf{Redis Sentinel} monitors a single primary-replica group with heartbeats; if a quorum of Sentinels detects \texttt{ODOWN}, they promote a replica automatically, though write throughput stays capped at one machine.

\textbf{Redis Cluster} shards writes horizontally across up to 1,000 nodes by dividing the keyspace into \textbf{16,384 hash slots}. Every key maps to a slot via:
\begin{equation}
\text{slot} = \text{CRC16}(k) \pmod{16384}
\end{equation}
Cluster nodes exchange slot ownership continuously via a binary \textbf{Gossip protocol}. If a client asks Master~1 for a key belonging to slot~12000 owned by Master~3, Master~1 replies \texttt{-MOVED 12000 10.0.0.4:6379} so the client updates its routing table. To force related keys onto the same node for multi-key operations or Lua scripts, wrap the shared identifier in \textbf{Hash Tags} (e.g., \texttt{\{user:8492\}:orders} and \texttt{\{user:8492\}:profile}): Redis hashes only the substring inside \texttt{\{...\}}.

\begin{figure}[htbp]
\centering
\begin{tikzpicture}[
  scale=0.84,
  box/.style={rectangle, draw=headingblue, fill=codebg, thick, rounded corners=2pt, minimum width=2.8cm, minimum height=1.2cm, align=center, font=\sffamily\tiny}
]
  \node[box, fill=exampleBg] (client) at (0, 0) {\textbf{Redis Client}\\Hashes \texttt{\{user:8492\}}\\$\text{CRC16}\bmod 16384 = 12000$};
  \node[box, fill=mistakeBg] (m1) at (5.2, 1.8) {\textbf{Master 1}\\Slots 0--5460};
  \node[box] (m2) at (5.2, 0.0) {\textbf{Master 2}\\Slots 5461--10922};
  \node[box, fill=intuitionBg] (m3) at (5.2, -1.8) {\textbf{Master 3}\\Slots 10923--16383};

  \draw[-{Stealth[scale=0.9]}, thick, color=warningRed] (client) -- node[above, sloped, font=\sffamily\tiny] {Wrong Node Query} (m1);
  \draw[-{Stealth[scale=0.9]}, thick, dashed, color=warningRed] (m1) to [bend right=20] node[above, sloped, font=\sffamily\tiny] {\texttt{-MOVED 12000 Master3}} (client);
  \draw[-{Stealth[scale=0.9]}, thick, color=intuitionTeal] (client) -- node[below, sloped, font=\sffamily\tiny] {Direct Retry (Slot 12000)} (m3);
  \draw[<->, dashed, color=gray!70] (m1) -- (m2);
  \draw[<->, dashed, color=gray!70] (m2) -- (m3);
\end{tikzpicture}
\caption{Redis Cluster 16,384 Hash-Slot Sharding and \texttt{-MOVED} Redirection}
\label{fig:v4-ch06-redis-cluster}
\end{figure}

\begin{realworldbox}[Diagnosing a Hot-Key Stampede and \texttt{CROSSSLOT} Error]
\textbf{Symptom:} Redis CPU hits 99\% on one shard while the application logs \texttt{ResponseError: CROSSSLOT Keys in request don't hash to the same slot}.\\
\textbf{Diagnostic Commands:}
\begin{verbatim}
$ redis-cli --hotkeys
# Top key: "homepage:banner" (412,000 req/s) -- TTL: 2s!
$ redis-cli INFO commandstats | grep cmdstat_mget
cmdstat_mget:calls=94000,rejected_calls=12400
\end{verbatim}
\textbf{Root Cause \& Fix:} (1) \texttt{homepage:banner} has a short TTL without XFetch, causing bursts of misses; apply XFetch or local in-process L1 caching. (2) Multi-key \texttt{MGET}/\texttt{MSET} or Lua calls across keys without matching \texttt{\{...\}} hash tags fail across Redis Cluster shards; align keys using hash tags or issue single-key pipelined calls.
\end{realworldbox}

\begin{tryityourselfbox}[Slot Math, XFetch Threshold, and \texttt{listpack} Sizing]
\begin{enumerate}[leftmargin=18pt]
  \item In \texttt{get\_with\_xfetch}, suppose \texttt{delta = 0.2s}, \texttt{beta = 1.0}, and \texttt{ttl\_rem = 0.4s}. What random value \texttt{u} must be drawn for $-\beta \cdot \delta \cdot \ln(u) > \text{ttl\_rem}$ to trigger an early refresh?
  \item Explain in two sentences why inserting a 300-byte field into a legacy Redis \texttt{ziplist} could take $O(N^2)$ time, while \texttt{listpack} stays $O(1)$ at the insertion tail.
\end{enumerate}
\end{tryityourselfbox}

\section{Summary \protect\& Key Takeaways}
\begin{itemize}[leftmargin=18pt]
  \item Redis pairs a single-threaded \texttt{epoll}/\texttt{kqueue} event loop with cache-friendly C structures (\texttt{sds}, incrementally rehashed \texttt{dict}, coin-flip \texttt{skiplist}, and \texttt{listpack} replacing \texttt{ziplist}).
  \item Eviction uses $N$-key sampled LRU and an 8-bit probabilistic Morris counter for LFU.
  \item Defend caches with TTL jitter (avalanche), XFetch early refresh (stampede), Bloom filters (penetration), and delete-on-write invalidation.
  \item Redis Cluster partitions keys across 16,384 CRC16 slots; multi-key commands and Lua scripts require matching \texttt{\{hash\_tags\}}.
\end{itemize}

\section{Exercises}

\subsection*{Exercise~\thechapter.1 (Mechanics): Redis Cluster Hash Tags}
Given keys \texttt{user:orders:8492} and \texttt{user:profile:8492}: (1) How does Redis Cluster locate each key? (2) Why does an atomic Lua script touching both fail with \texttt{CROSSSLOT}? (3) Rewrite both keys with hash tags so they map to the same slot.

\subsection*{Exercise~\thechapter.2 (Debugging): Stale Email Overwrite Race}
Identify the concurrency race in the code below and provide the production fix:
\begin{lstlisting}[language=Python]
def update_user_email(user_id, new_email):
    db.execute("UPDATE users SET email = %s WHERE id = %s", (new_email, user_id))
    redis.set(f"user:{user_id}", new_email)
\end{lstlisting}

\subsection*{Exercise~\thechapter.3 (Architecture): Global Distributed Rate Limiter}
Design a Redis rate limiter enforcing 1,000 requests/minute per API key across 200 frontend servers. Compare a \texttt{ZSET} sliding-window log vs. a Lua Token Bucket using \texttt{HSET}, and write the Lua script.

\subsection*{Exercise~\thechapter.4 (Try-It-Yourself Lab): XFetch Probability & \texttt{listpack}}
Solve both parts of the \emph{Try It Yourself} box above.

\section{Detailed Solutions}

\subsection*{Solution to Exercise~\thechapter.1}
\begin{enumerate}[leftmargin=18pt]
  \item Redis computes $\text{CRC16}(\text{key}) \pmod{16384}$ and routes to the master owning that slot.
  \item Without hash tags, the two strings hash to different slots on different cluster masters; Redis Cluster forbids multi-key operations across slots.
  \item Rewrite as \texttt{\{user:8492\}:orders} and \texttt{\{user:8492\}:profile}; both hash strictly on \texttt{"user:8492"}.
\end{enumerate}

\subsection*{Solution to Exercise~\thechapter.2}
If Thread~1 writes Email~1 to the DB, pauses, Thread~2 writes Email~2 to DB and sets Redis to Email~2, and then Thread~1 resumes and calls \texttt{redis.set()}, Redis is overwritten with stale Email~1. Fix by deleting the key after DB commit: \texttt{redis.delete(f"user:\{user\_id\}")}.

\subsection*{Solution to Exercise~\thechapter.3}
A \texttt{ZSET} log stores every request timestamp ($O(N)$ RAM per user), whereas a Lua Token Bucket stores just two hash fields (\texttt{tokens}, \texttt{last\_updated}, $O(1)$ RAM). Note that \texttt{HSET} replaces deprecated \texttt{HMSET}:
\begin{lstlisting}[language=Python]
LUA_RATE_LIMITER = """
local key = KEYS[1]
local capacity = tonumber(ARGV[1])
local refill_rate = tonumber(ARGV[2])
local now = tonumber(ARGV[3])
local requested = tonumber(ARGV[4])

local data = redis.call('HMGET', key, 'tokens', 'last_updated')
local tokens = tonumber(data[1])
local last_updated = tonumber(data[2])
if not tokens then
    tokens = capacity
    last_updated = now
else
    local delta = math.max(0, now - last_updated)
    tokens = math.min(capacity, tokens + delta * refill_rate)
end

if tokens >= requested then
    tokens = tokens - requested
    redis.call('HSET', key, 'tokens', tokens, 'last_updated', now)
    redis.call('EXPIRE', key, math.ceil(capacity / refill_rate))
    return 1
else
    return 0
end
"""
\end{lstlisting}

\subsection*{Solution to Exercise~\thechapter.4}
\begin{enumerate}[leftmargin=18pt]
  \item We need $-0.2 \ln(u) > 0.4 \implies \ln(u) < -2 \implies u < e^{-2} \approx 0.135$ (a 13.5\% chance per read when $0.4\text{ s}$ remains).
  \item In \texttt{ziplist}, each entry's \texttt{prevlen} field grew from 1 to 5 bytes whenever the preceding entry crossed 254 bytes, forcing $O(N^2)$ cascading reallocations across all following entries. In \texttt{listpack}, each entry stores only its own length at its tail, decoupling neighbors so updates never cascade.
\end{enumerate}

\begin{labbridgebox}[Redis Caching & Idempotency in Production]
Inspect live Redis~7 caching and Lua/key invalidation alongside PostgreSQL and RabbitMQ inside \texttt{projects/p12-end-to-end-platform} (port 6379).
\end{labbridgebox}

\section{Further Reading \protect\& Citations}
\begin{itemize}[leftmargin=18pt]
  \item Sanfilippo, S. (2020). \emph{Redis Documentation and Source Code}. \url{https://redis.io}.
  \item Vattani, A., Chierichetti, F., \& Lowenstein, K. (2015). \emph{Optimal probabilistic cache stampede prevention}. VLDB Endowment, 8(8), 886--897.
  \item Pugh, W. (1990). \emph{Skip lists: a probabilistic alternative to balanced trees}. CACM, 33(6), 668--676.
\end{itemize}
